% =============================================================================
% Kapitel 8: Diskussion, Reflexion und Beantwortung der Forschungsfragen
% ARBEITSGERUEST / REVIEW CANDIDATE
% Stand: 03.09.2026
%
% Dieses Dokument enthaelt bewusst NOCH KEINEN ausformulierten Thesis-Text.
%
% Kapitel 7 = Was wurde beobachtet?
% Kapitel 8 = Was bedeuten die beobachteten Ergebnisse?
%
% ZENTRALE AUFGABEN VON KAPITEL 8:
%
% 1. Ergebnisse aus Kapitel 7 interpretieren.
% 2. Ergebnisse mit Literatur und Stand der Technik abgleichen.
% 3. Spannungen, Widersprueche und Grenzen offen diskutieren.
% 4. technische Machbarkeit von wissenschaftlicher / fachlicher Validitaet trennen.
% 5. Effizienz, Kosten, Laufzeit und Human Interaction reflektieren.
% 6. Safe-Demo-Feedback wissenschaftlich einordnen.
% 7. AI im Turing Generator und AI als Entwicklungswerkzeug getrennt reflektieren.
% 8. HF / TF1 / TF2 / TF3 explizit beantworten.
%
% HARTER SCHREIBGRUNDSATZ:
%
% Kapitel 8 darf interpretieren, aber nur auf Basis von:
% - Ergebnissen aus Kapitel 7,
% - wissenschaftlicher Literatur,
% - explizit dokumentierter Projektevidenz.
%
% KEINE neue empirische Beobachtung in Kapitel 8 einfuehren,
% die nicht zuvor in Kapitel 7 dokumentiert wurde.
%
% QUELLENLOGIK:
%
% wissenschaftlicher General Claim
% -> wissenschaftliche Literatur
%
% projektspezifische Beobachtung
% -> Kapitel 7 / Projektartefakt
%
% Interpretation der Beobachtung
% -> Kombination aus Ergebnis + Literatur
%
% ADR
% -> Design Rationale / Projektevidenz, NICHT Literatur.
%
% MUST:
% - keine Forschungsfrage stillschweigend umformulieren
% - keine universellen Aussagen aus einem PoC ableiten
% - keine Ontologie als alleinige Anti-Hallucination-Loesung darstellen
% - SYSIDE-Validitaet nicht mit Engineering Correctness gleichsetzen
% - Testanzahl nicht als wissenschaftliche Validitaet interpretieren
% - Multi-Source-Scope nicht ueber Project 308131 hinaus verallgemeinern
%
% =============================================================================

\section{Diskussion, Reflexion und Beantwortung der Forschungsfragen}
\label{sec:discussion}

% =============================================================================
\subsection{Diskussionsrahmen}
\label{sec:discussion_frame}
% =============================================================================

\subsubsection{Bezug zwischen Ergebnissen, Literatur und Forschungsfragen}
\label{sec:discussion_evidence_logic}

% ZWECK:
% Leser kurz erklaeren, wie Kapitel 8 aufgebaut ist.
%
% LOGIK:
%
% Ergebnis aus Kapitel 7
% -> Interpretation
% -> Abgleich mit Literatur
% -> Limitation / Gegenbefund
% -> Beitrag zur Forschungsfrage
%
% KEINE erneute ausfuehrliche Ergebniswiedergabe.
%
% Beispielstruktur pro Diskussionsblock:
%
% 1. 1-2 Saetze Ergebnisreferenz
% 2. Interpretation
% 3. Literaturabgleich
% 4. Claim Boundary
% 5. RQ-Relevanz
%
% NICHT:
% Kapitel 7 paraphrasieren.


\subsubsection{Ebenen der Diskussion}
\label{sec:discussion_levels}

% Die Diskussion unterscheidet vier Ebenen:
%
% A. Architektur:
%    Welche Verantwortungstrennungen waren tragfaehig?
%
% B. KI-gestuetzte Informationsverarbeitung:
%    Wo ist LLM-Unterstuetzung sinnvoll und wo braucht es deterministische /
%    Human-controlled Boundaries?
%
% C. Modellbildung / SysML v2:
%    Welche Trennung ist zwischen Meaning, Representation, Formulation,
%    Assembly und Serialization erforderlich?
%
% D. Untersuchung / PoC:
%    Wie belastbar und uebertragbar sind die Ergebnisse?
%
% Dadurch nicht Architektur-, Tool- und Evaluationserkenntnisse vermischen.


% =============================================================================
\subsection{Interpretation der Ergebnisse zur source-grounded Informationsverarbeitung}
\label{sec:discussion_source_grounding}
% =============================================================================

\subsubsection{Source Purity als notwendige Voraussetzung}
\label{sec:discussion_source_purity}

% ERGEBNISBASIS:
% Kapitel 7:
% - SEM-001
% - BLK-003
% - fruehe Processing Runs vs R4c
%
% ZU INTERPRETIEREN:
%
% 1. Warum ist die Trennung von Engineering Source und Processing Context
%    architektonisch relevant?
%
% 2. Welche Fehler entstehen, wenn Prompt-/Task-/Orchestration-Information
%    als fachliche Evidence behandelt wird?
%
% 3. Was bedeutet das fuer TF1:
%    Nicht nur Dokumentqualitaet, sondern kontrollierter Source Context /
%    Evidence Boundary ist Voraussetzung.
%
% MOEGLICHE KERNAUSSAGE:
%
% Zuverlaessige KI-Verarbeitung setzt nicht nur "gute Eingangsdaten" voraus.
% Ebenso wichtig ist eine explizite maschinelle Grenze zwischen:
%
% Engineering Source Content
% und
% Processing / Reference / Orchestration Context.
%
% LITERATURABGLEICH:
%
% [REPO-BIB]
% Chis 2025 / heterogeneous KB / RAG
% Tian 2025 / heterogeneous multi-source data
% ggf. provenance / context literature
%
% [ADD / PRUEFEN]
% LLM grounding / context contamination / provenance source
%
% CLAIM BOUNDARY:
% PoC-Evidenz zeigt dies im untersuchten Workflow,
% nicht fuer alle LLM-basierten Engineering-Systeme.


\subsubsection{Evidence und Subject Identity vor Persona Interpretation}
\label{sec:discussion_subject_identity}

% ERGEBNISBASIS:
% SEM-003
% SEM-004
% BLK-003
% R4c-Recovery
%
% ZU INTERPRETIEREN:
%
% Frueher:
% Persona erzeugt Proposal Population
% -> nachgelagerte Consolidation
%
% Beobachtet:
% D3 109 -> D4 109
%
% Spaeter:
% source-grounded Evidence
% -> shared Subject Identity
% -> Persona Interpretation
%
% DISKUSSIONSFRAGE:
%
% Warum ist eine gemeinsame Identitaet vor Multi-Persona-Verarbeitung
% methodisch / architektonisch robuster als spaetere Rekonstruktion?
%
% WICHTIG:
% Nicht allein aus 109->109 universell schliessen.
% Kombination:
% - empirische Beobachtung
% - weitere semantische Findings
% - Architekturkorrektur
% - Retest
%
% POTENZIELLER THESIS-BEITRAG:
%
% Persona Diversity sollte auf denselben source-grounded Gegenstand
% angewendet werden.
%
% Persona diversity != unabhängige Subject discovery.
%
% LITERATUR:
% [ADD]
% Multi-agent disagreement / diversity:
% Wang 2023
% Liang 2024
% Du 2024
% Cui 2026
% Kostka & Chudziak 2026
%
% Vor Verwendung Volltexte claimweise pruefen.


\subsubsection{Review-Item Count als ungeeignetes Primaerziel}
\label{sec:discussion_review_count}

% ERGEBNISBASIS:
% SEM-005
%
% ZU INTERPRETIEREN:
%
% Weniger Review Items koennen:
% - bessere Konsolidierung bedeuten
% - ODER Informationsverlust bedeuten
%
% Mehr Review Items koennen:
% - echte Informationsdichte darstellen
% - ODER kuenstliche Persona-/Run-Multiplikation
%
% Deshalb:
% Count alleine keine sinnvolle Assurance-/Qualitaetsmetrik.
%
% Kandidaten fuer geeignetere Kriterien:
% - Source Purity
% - Grounding
% - Relevance
% - Coverage
% - echte Variance
% - Human effort
% - unresolved ambiguity
%
% WICHTIG:
% Diese Kriterien als Interpretation diskutieren,
% nicht als bereits statistisch validierte Metrik verkaufen.


% =============================================================================
\subsection{Interpretation der Ergebnisse zu Unsicherheit, Persona-Varianz und Human Authority}
\label{sec:discussion_human_authority}
% =============================================================================

\subsubsection{Legitime Unsicherheit als Engineering Information}
\label{sec:discussion_uncertainty}

% ERGEBNISBASIS:
% SEM-008
% PASS-004
% Open Questions in Project 120412
%
% ZENTRALE INTERPRETATION:
%
% Unsicherheit muss nicht automatisch:
% - erzwungen aufgeloest,
% - halluzinatorisch ergaenzt,
% - oder als technischer Failure behandelt werden.
%
% Sinnvolle Trennung:
%
% semantic uncertainty
% -> Human Engineering Review
%
% integrity / provenance / identity violation
% -> fail closed
%
% DISKUSSION:
% Das ist eine wichtige Assurance-Grenze.
%
% LITERATUR:
% [ADD]
% HITL / Human oversight:
% Lazaros 2026
% Mosqueira-Rey 2023
% Zhu 2026
%
% [REPO-BIB]
% ggf. GRADE nur wenn Claim passt.
%
% TF2-Relevanz hoch.


\subsubsection{Consensus, Shared Ambiguity und Persona Variance}
\label{sec:discussion_variance}

% ERGEBNISBASIS:
% SEM-013
% OBS-031
%
% Drei unterschiedliche Zustandsarten:
%
% 1. consensus on one interpretation
% 2. shared ambiguity
% 3. actual Persona variance
%
% INTERPRETATION:
%
% Diese drei Zustaende sollten nicht dieselbe Human-Interaction ausloesen.
%
% Effective Intervention:
%
% single consensus:
% -> geringe Interaktionslast
%
% shared ambiguity:
% -> Human waehlt / klaert gemeinsame Alternativen
%
% actual variance:
% -> Differenzen explizit untersuchen
%
% WICHTIG:
% Human Approval darf bei Consensus trotzdem nicht implizit werden,
% wenn Engineering Authority erforderlich ist.
%
% Literatur:
% Multi-agent disagreement / ensemble / confidence
% nach Volltextpruefung.
%
% TF2-Kernbereich.


\subsubsection{Human Authority als gestufte statt einmalige Kontrolle}
\label{sec:discussion_staged_authority}

% ERGEBNISBASIS:
%
% Kapitel 7:
% - Engineering Review
% - Model Placement Review
% - Model Quality / Formulation Review
% - Assembly relationship resolution
% - Final Model Review
% - Publication Approval
%
% INTERPRETATION:
%
% "Human in the Loop" ist fuer diesen Use Case kein einzelner
% generischer Approve-Schritt.
%
% Unterschiedliche Entscheidungstypen:
%
% Meaning Authority
% Placement Authority
% Formulation / Quality Authority
% Model Authority
% Publication Authority
%
% Diese Entscheidungen haben unterschiedliche Informationsbedarfe.
%
% TF2:
% Effective Intervention bedeutet nicht maximale Human Interaction.
% Es bedeutet Human Authority an materiellen Engineering-Entscheidungen,
% kombiniert mit Reduktion unnoetiger Interaktion.
%
% Literatur:
% HITL / Human oversight / meaningful human control.
%
% Wichtig:
% nicht behaupten, dass exakt diese Gate-Anzahl universell optimal ist.


\subsubsection{Revidierbarkeit und Auditierbarkeit menschlicher Entscheidungen}
\label{sec:discussion_human_revision}

% ERGEBNISBASIS:
% OBS-020
% OBS-021
% OBS-025
% OBS-026
%
% INTERPRETATION:
%
% Human Authority ist nicht nur "jemand klickt Approve".
%
% Eine belastbare Authority Chain benoetigt:
%
% - eindeutige Action Semantics
% - Reviewer Binding
% - immutable Decision Revision
% - explizites Reopen / Supersede
% - nachvollziehbaren finalen Authority State
%
% Literatur:
% Provenance / auditability / human oversight.
%
% W3C PROV ggf. fuer provenance,
% aber nicht als Human-Authority-Norm ueberdehnen.


% =============================================================================
\subsection{Interpretation der Ergebnisse zur Modellableitung}
\label{sec:discussion_model_derivation}
% =============================================================================

\subsubsection{Engineering Meaning ist nicht Target Model Representation}
\label{sec:discussion_meaning_representation}

% ERGEBNISBASIS:
% SEM-002
% SEM-012
% BLK-006
%
% ZENTRALE INTERPRETATION:
%
% Engineering Information Classification:
% beschreibt Bedeutung der extrahierten Information.
%
% Target Model Representation:
% beschreibt, wie diese Information im Zielmodell repraesentiert wird.
%
% Beides ist nicht zwingend 1:1.
%
% Beispiel aus Projekt:
% Engineering Information = constraint
% Target Representation = stakeholder requirement
%
% DISKUSSION:
% Warum ist ein direktes Mapping problematisch?
% Welche Verantwortung muss zwischen beiden Ebenen liegen?
%
% Literaturabgleich:
% - Model Transformation / Metamodelle
% - QVT / MDA
% - SysML transformation approaches
%
% [REPO-BIB]
% Kapos 2014
% Li 2022
% Zhou 2024
% Meyers 2013
% Zhao 2023
%
% TF3 sehr relevant.


\subsubsection{Placement, Formulation und Serialization als getrennte Verantwortungen}
\label{sec:discussion_formulation}

% ERGEBNISBASIS:
% BLK-006
% SEM-015
% SEM-014
%
% Finales Konzept:
%
% Engineering Meaning
% -> Model Placement
% -> Target-Model Formulation
% -> deterministic Assembly
% -> SysML v2 Serialization
%
% INTERPRETATION:
%
% Placement beantwortet:
% Wo / als welche Modellrolle soll Information repraesentiert werden?
%
% Formulation beantwortet:
% Wie wird akzeptierte Bedeutung modellgerecht formuliert?
%
% Serialization beantwortet:
% Wie wird ein autorisiertes Modell deterministisch in Syntax ueberfuehrt?
%
% Diese Trennung verhindert, dass Syntaxgenerator implizit Engineering Meaning
% erfinden oder umdeuten muss.
%
% LITERATUR:
% SysML v2 spec
% model transformation
% LLM SysML generation papers
% Architecture as Code
%
% [REPO-BIB]
% Cibrian 2025
% Cibrian 2025b
% Pan 2025
% Ahlbrecht 2024
% SysMLv2
%
% [ADD]
% July 2026 Prompt-Strategy-Driven SysML-v2 Artefact Generation
%
% WICHTIG:
% "verhindert Halluzination" nicht absolut schreiben.
% Besser:
% reduziert den Entscheidungsspielraum downstream / macht Entscheidungen explizit.


\subsubsection{Deterministische Assembly und Generation}
\label{sec:discussion_deterministic_generation}

% ERGEBNISBASIS:
% BLK-007
% IEM
% Generation / SYSIDE
%
% INTERPRETATION:
%
% Nach Human-authorized Meaning + Placement:
% keine erneute freie LLM-Interpretation in Assembly / Serialization.
%
% Vorteile im PoC:
% - reproduzierbarer downstream artifact path
% - explizite unsupported semantics
% - fail-closed behavior
% - klare Validierungsgrenze
%
% Tradeoff:
% - begrenzte Construct Coverage
% - neue Constructs brauchen explizite Profiles / Mappings
%
% Literatur:
% deterministic transformation
% metamodel validation
% SysML v2 generation.
%
% TF3 / HF.


\subsubsection{Formale Modellvaliditaet und Engineering Correctness}
\label{sec:discussion_validity_correctness}

% ERGEBNISBASIS:
% SYSIDE PASS
% Human Final Model Review
%
% ZENTRALE TRENNUNG:
%
% external tool validation
% !=
% Engineering correctness
%
% SYSIDE prueft:
% akzeptierte syntax/tool compatibility im verwendeten Scope.
%
% Human Final Review prueft:
% fachliche Plausibilitaet / autorisierte Modellentscheidung im PoC.
%
% Keines ersetzt das andere.
%
% Literatur:
% SysML v2 spec
% Cibrian metamodel validation
% Verification/Validation sources.
%
% Diese Trennung unbedingt in HF/TF3-Antwort aufnehmen.


% =============================================================================
\subsection{Interpretation der Multi-Source-Ergebnisse}
\label{sec:discussion_multisource}
% =============================================================================

\subsubsection{Project-wide Identity und Provenance}
\label{sec:discussion_multisource_identity}

% ERGEBNISBASIS:
% BLK-002 initial FAIL
% Project 308131 Retest PASS
%
% INTERPRETATION:
%
% Source-local Eindeutigkeit reicht nicht,
% sobald mehrere Sources in einen Project-level Prozess eingehen.
%
% Benoetigt:
% project-aware identity
% source/run/artifact provenance
% exact admission
%
% Diskussion mit:
% W3C PROV
% multi-source data integration literature
% heterogeneous engineering information literature
%
% [ADD]
% W3C PROV-DM / PROV Primer
% optional Simmhan 2005
%
% TF1 / HF.


\subsubsection{Source Neutrality und Konfliktbehandlung}
\label{sec:discussion_source_neutrality}

% ERGEBNISBASIS:
% Project 308131 acceptance semantics
%
% Keine a-priori Source Authority Hierarchy.
%
% Equivalent:
% kann zusammengefuehrt / gemeinsam weitergefuehrt werden,
% Provenance bleibt erhalten.
%
% Distinct / complementary:
% koexistiert.
%
% Conflict / uncertainty:
% bleibt explizit / Human Authority.
%
% INTERPRETATION:
%
% "neueste Source gewinnt" oder "Source Type X ist autoritativ"
% wurde bewusst nicht generalisiert.
%
% Literatur:
% heterogeneous information integration
% provenance
% data fusion / conflict resolution, falls passende Literatur vorhanden.
%
% WICHTIG:
% Nur accepted Project-308131-Scope diskutieren.
% Automatic cross-source change control bleibt Outlook.


\subsubsection{Grenzen des Genuine-Multi-Source-Nachweises}
\label{sec:discussion_multisource_limits}

% Klarstellen:
%
% Project 308131:
% genuine Multi-Source PoC evidence
%
% NICHT nachgewiesen:
% - beliebig viele Sources
% - langfristiges Change Management
% - automatische Authority Resolution
% - industrielle Scale
% - beliebige Source Types
%
% Concern-centric S3A/S3B/S4/S5:
% prototype evidence only.
%
% Diese Grenzen spaeter auch Threats to Validity zuordnen.


% =============================================================================
\subsection{Abgleich mit Stand der Technik und Forschungsluecke}
\label{sec:discussion_literature_gap}
% =============================================================================

\subsubsection{LLM-basierte SysML-v2-Generierung}
\label{sec:discussion_llm_sysml}

% Literatur aus Kapitel 2 wieder aufnehmen:
%
% Pan 2025
% Cibrian 2025
% Cibrian 2025b
% Prompt-Strategy-Driven SysML-v2 Artefact Generation
%
% VERGLEICHSACHSEN:
%
% - Eingangsquelle / Heterogenitaet
% - Source Grounding
% - Human Authority
% - Multi-Source
% - Provenance
% - Model Placement
% - Target Formulation
% - deterministic generation
% - external validation
%
% Nicht paper-by-paper referieren.
% Besser Matrix / thematischer Abgleich.
%
% Forschungsbeitrag nicht als "erstes System ueberhaupt" behaupten.
% Formulierung:
% "In der durchgefuehrten Literaturrecherche wurde kein Ansatz identifiziert,
% der die untersuchten Verantwortungen in dieser Kombination abbildet."
%
% Nur wenn LSR dies nach finalem Screening stuetzt.


\subsubsection{Ontologien und semantische Leitplanken}
\label{sec:discussion_ontology_guardrails}

% TF3 fragt explizit nach ontologischen Regeln / Leitplanken.
%
% Ergebnis des Projekts duerfte differenzierter sein:
%
% Ontologie / semantic profiles sind hilfreich fuer:
% - Terminologie
% - Klassifikation
% - erlaubte Beziehungen / types
% - controlled semantics
%
% Aber im PoC waren zusaetzlich erforderlich:
%
% - Source Grounding
% - Provenance
% - Identity
% - Human Authority
% - Project Fit
% - explicit model placement
% - Target-Model Formulation
% - deterministic assembly
% - validation
%
% ZENTRALE DISKUSSION:
%
% TF3-Annahme ggf. praezisieren:
% Ontologische Regeln allein sind keine ausreichende Architektur.
%
% WICHTIG:
% Forschungsfrage nicht umschreiben.
% Antwort darf zeigen, dass die Ausgangsannahme zu eng war.
%
% Literatur:
% Dong
% Drobnjakovic / IOF
% ISO/IEC 21838-2
% Liu SysDO
% Lu
% Orellana/Mandrick
%
% Keine Aussage "Ontology prevents hallucination".


\subsubsection{Architecture as Code}
\label{sec:discussion_architecture_as_code}

% Abgleich:
%
% Architecture as Code:
% maschinenlesbare / versionierbare Architekturartefakte
%
% Turing Generator:
% generiert textual SysML v2 als kontrolliertes Zielartefakt.
%
% Diskussion:
% - versionierbarkeit
% - deterministic generation
% - traceability
% - reproducibility
% - human approval
%
% Literatur:
% Bucaioni et al. ICSA 2025
% Stirbu et al. IEEE Software 2022
% ggf. Pontillo ICSA 2026
%
% Nicht behaupten, dass jede SysML-v2-Textdatei automatisch Architecture as Code ist.
% Bedingungen explizieren.


% =============================================================================
\subsection{Effizienz, Laufzeit und praktische Nutzbarkeit}
\label{sec:discussion_efficiency}
% =============================================================================

\subsubsection{LLM-Aufrufe, Laufzeit und Kosten}
\label{sec:discussion_runtime_cost}

% ERGEBNISBASIS:
% SEM-007
% OBS-017
% ggf. spaetere Token-/Runtime-Metriken
%
% NOCH ZU ERMITTELN / ZU SAMMELN:
%
% - LLM Call Count representative Runs
% - Input / Output Tokens, falls Logging vorhanden
% - Laufzeiten representative Processing Steps
% - ggf. API-Kosten fuer representative Runs
%
% KEINE Zahlen erfinden.
%
% DISKUSSION:
%
% Multi-Persona / Repetition verbessert nicht automatisch den Nutzen,
% verursacht aber Calls / Laufzeit / Review-Volumen.
%
% R4c reduziert unnoetige Proposal Multiplication,
% aber quantitative Effizienzverbesserung nur behaupten,
% wenn Daten vorhanden.
%
% Kapitel 8 darf qualitative Tradeoffs diskutieren,
% quantitative Claims brauchen Messdaten.


\subsubsection{Human Interaction und Automatisierungsgrad}
\label{sec:discussion_human_efficiency}

% ERGEBNISBASIS:
% SEM-005
% SEM-013
% OBS-031
% Human Review Gates
%
% SPANNUNG:
%
% mehr Human Gates:
% + kontrollierte Engineering Authority
% - mehr Interaktionsaufwand
%
% Ziel:
% Authority nicht minimieren,
% sondern Human Attention auf materielle Entscheidungen konzentrieren.
%
% Diskutieren:
% - consensus one-click
% - shared ambiguity targeted choice
% - true variance deep review
% - unimportant technical metadata progressive disclosure
%
% TF2-Relevanz hoch.


\subsubsection{Praxistauglichkeit des PoC}
\label{sec:discussion_practicality}

% ERGEBNISBASIS:
% E2E
% Runtime
% Streamlit coupling
% Safe Demo spaeter
%
% Diskutieren:
%
% PoC demonstrierbar / technisch end-to-end
%
% aber nicht Production Ready:
% - Background execution
% - cancellation
% - reprocessing / supersession
% - broader target coverage
% - performance
% - usability
%
% Keine Produktreife suggerieren.


% =============================================================================
\subsection{Einordnung der externen Safe Demo}
\label{sec:discussion_safe_demo}
% =============================================================================

% STATUS:
% erst nach Safe Demo ausformulieren.


\subsubsection{Abgleich mit den internen Evaluationsergebnissen}
\label{sec:discussion_safe_demo_alignment}

% NACH DER DEMO:
%
% Welche internen Beobachtungen wurden durch Practitioner Feedback:
% - bestaetigt?
% - relativiert?
% - nicht angesprochen?
% - neu ergaenzt?
%
% Wichtig:
% keine "Validation by experts" behaupten, wenn es nur Demo + Feedback war.


\subsubsection{Aussagekraft der Practitioner-Rueckmeldung}
\label{sec:discussion_safe_demo_validity}

% Grenzen:
% - kleine / nicht representative Gruppe
% - Demonstrationskontext
% - keine kontrollierte Vergleichsgruppe
% - moegliche Gruppendynamik
% - Entwickler praesentiert eigenes System
%
% Staerke:
% - externe fachliche Perspektive
% - Praxisplausibilisierung
% - konkrete Fragen / Bedenken / Akzeptanzsignale
%
% Nur tatsaechlich dokumentiertes Feedback verwenden.


% =============================================================================
\subsection{Reflexion des KI-gestuetzten Entwicklungsprozesses}
\label{sec:discussion_ai_development}
% =============================================================================

% WICHTIGE TRENNUNG:
%
% AI ROLE A:
% KI als Bestandteil des untersuchten Turing Generators
%
% AI ROLE B:
% KI als Entwicklungswerkzeug / Coding- und Architecture Assistant
%
% Nicht vermischen.


\subsubsection{KI als Bestandteil des untersuchten Systems}
\label{sec:discussion_ai_system_role}

% Zusammenfassen:
%
% AI beteiligt an:
% - Evidence / semantic interpretation, soweit finaler Flow
% - Persona Interpretation
% - Model Derivation / Formulation in bounded stages
%
% Keine Engineering Approval Authority.
%
% Diskutieren:
% wo probabilistische Interpretation sinnvoll war
% und wo deterministische / Human-controlled Verarbeitung bevorzugt wurde.


\subsubsection{KI als Entwicklungs- und Coding-Werkzeug}
\label{sec:discussion_ai_coding}

% PROJEKTSPEZIFISCHE REFLEXION:
%
% Entwicklungsworkflow:
%
% AI proposal
% -> Human Review
% -> ADR / Architecture Decision
% -> bounded slice
% -> automated tests
% -> observed failure
% -> root-cause analysis
% -> architecture correction
% -> retest / regression
% -> Human acceptance
%
% Potenzielle Beobachtungen:
%
% - Code-Erzeugung wurde stark beschleunigt.
% - Human-Arbeit verschob sich teilweise von Zeilenimplementierung zu:
%   Architecture Boundaries
%   Contracts
%   Test Interpretation
%   Root-Cause Analysis
%   Scope
%   Acceptance
%
% ABER:
% Diese Aussagen nur als Projekterfahrung formulieren,
% sofern keine quantitativen Daten vorliegen.
%
% NICHT:
% allgemeine Produktivitaetssteigerung fuer Softwareentwicklung behaupten.
%
% Literatur:
% Agentic Systems as Boosting Weak Reasoning Models
% Code as Agent Harness
% ggf. wissenschaftliche AI-assisted development studies
%
% [ADD / PRUEFEN]
% aktuelle empirische Literatur zu AI coding assistants.


\subsubsection{Human Authority auch im Entwicklungsprozess}
\label{sec:discussion_ai_development_authority}

% Interessanter Meta-Befund:
%
% Dieselbe Grundidee wie im Turing Generator:
% AI kann Vorschlaege erzeugen,
% Engineering Authority bleibt beim Menschen.
%
% ABER:
% Vorsicht vor zu eleganter Selbstaehnlichkeit.
% Nicht behaupten, dies sei empirisch "bewiesen".
%
% Diskutierbar:
% - Architekturfreigabe
% - ADR Acceptance
% - Test Interpretation
% - Scope Freeze
% - finaler CATIA Model Freeze
%
% Human war Entscheidungstraeger.


\subsubsection{Risiken des KI-gestuetzten Entwicklungsprozesses}
\label{sec:discussion_ai_development_risks}

% Diskutieren:
%
% - plausible, aber falsche Implementierungsvorschlaege
% - lokale Fixes, die Architekturgrenzen verletzen
% - Test Overfitting
% - Kontextverlust bei langen Entwicklungsverlaeufen
% - Gefahr, Tests als Ersatz fuer Engineering Review zu behandeln
% - hoher Reviewbedarf bei umfangreichen Codeaenderungen
%
% Gegenmassnahmen im Projekt:
%
% - ADRs
% - SSOT Checkpoints
% - bounded slices
% - focused tests
% - complete regression
% - Human Acceptance
% - CATIA architecture authority
%
% Generalisierung nur mit Literatur.


% =============================================================================
\subsection{Limitationen und Threats to Validity}
\label{sec:discussion_limitations}
% =============================================================================

\subsubsection{Formative Selbstevaluation und Entwicklerbias}
\label{sec:discussion_self_evaluation_limit}

% WP-12 Evaluator = Entwickler / Systems Engineer.
%
% Konsequenzen:
% - nicht unabhaengig
% - Erwartungsbias moeglich
% - qualitative Beobachtung
%
% Mitigation:
% - vorab definierte Protokolle
% - fail states erhalten
% - Findings Register
% - immutable checkpoints
% - Safe Demo als externe Zusatzperspektive
%
% Keine Behauptung, dass Bias eliminiert wurde.


\subsubsection{Veraenderung des Systems waehrend formativer Evaluation}
\label{sec:discussion_formative_change_limit}

% WP-12 war formative Evaluation:
% System wurde zwischen Findings korrigiert.
%
% Vorteil:
% Artifact Improvement / Architecture Learning
%
% Limitation:
% kein fixed-build Benchmark.
%
% Deshalb:
% formative Evidenz und final acceptance getrennt.


\subsubsection{Begrenzte Eingangs- und Domaenenabdeckung}
\label{sec:discussion_input_scope_limit}

% Eingangsquellen:
% bounded demonstrator / controlled data
%
% Genuine Multi-Source:
% Project 308131
%
% NICHT:
% industrielle Datenmenge / alle Legacy-Typen.
%
% TF1-Antwort deshalb als Constraints im untersuchten Scope formulieren.


\subsubsection{LLM-Variabilitaet und Reproduzierbarkeit}
\label{sec:discussion_llm_variability_limit}

% LLM-Ausgaben koennen variieren.
%
% Projekt-Gegenmassnahmen:
% - pinned configuration, soweit vorhanden
% - immutable processing artifacts
% - fingerprints
% - provenance
% - Human decisions
% - repeated runs als stability evidence, NICHT votes
%
% Aber:
% semantische Outputs nicht vollstaendig deterministisch.
%
% Keine uebertriebene Reproduzierbarkeitsbehauptung.


\subsubsection{Target-Model-Coverage}
\label{sec:discussion_target_scope_limit}

% SEM-011
% SEM-015 partial
%
% Nicht alle SysML-v2-Constructs abgedeckt.
%
% Daher:
% keine allgemeine SysML-v2-Generierungsfaehigkeit behaupten.
%
% Architektur kann erweiterbar konzipiert sein,
% Extension selbst aber nicht vollstaendig validiert.


\subsubsection{Tool- und Validierungsgrenzen}
\label{sec:discussion_tool_limit}

% CATIA:
% Architekturmodellierungswerkzeug / authority for Turing architecture.
%
% SYSIDE:
% externe Target Artifact Validation.
%
% SysML v2:
% official OMG specification.
%
% Grenzen:
% - Tool-Versionen
% - locally validated syntax patterns
% - keine Garantie anderer Toolchains
%
% Wichtig:
% CATIA model ist NICHT Syntax Authority fuer generierte Zielmodelle.


\subsubsection{Usability und Production Readiness}
\label{sec:discussion_usability_limit}

% Keine unabhaengige Usability Study.
%
% Keine statistische UX-Auswertung.
%
% Streamlit PoC.
%
% Runtime-/background-/cancel-/reprocessing findings.
%
% Safe Demo != usability validation.
%
% Production Readiness explizit ausserhalb des Claims.


% =============================================================================
\subsection{Uebertragbarkeit}
\label{sec:discussion_transferability}
% =============================================================================

\subsubsection{Uebertragbare Architekturprinzipien}
\label{sec:discussion_transferable_principles}

% Nicht "das System ist universell uebertragbar".
%
% Stattdessen:
% Welche Prinzipien KOENNTEN bei aehnlichen Problemklassen relevant sein?
%
% Kandidaten:
% - source-grounded Evidence
% - explicit provenance
% - identity before multi-agent interpretation
% - uncertainty preservation
% - staged Human Authority
% - separation Meaning / Representation / Formulation / Serialization
% - deterministic downstream generation
% - independent validation
%
% Generalisierung mit Literatur abgleichen.


\subsubsection{Voraussetzungen fuer andere Brownfield-Szenarien}
\label{sec:discussion_transfer_conditions}

% Conditions:
%
% - Source muss technisch zugreifbar / projizierbar sein
% - ausreichende inhaltliche Evidenz
% - kontrollierter Project Context
% - definierte Target Profile
% - Human Domain Expertise
% - externe Validierungsmethode
%
% Unterscheiden:
% Architekturprinzip uebertragbar
% vs konkrete Taxonomie / Prompts / Profiles projektspezifisch.


\subsubsection{Grenzen der Uebertragbarkeit}
\label{sec:discussion_transfer_limits}

% Unbekannt / nicht validiert:
%
% - andere Industriesektoren
% - andere Modellierungssprachen
% - sehr grosse Repositories
% - multimodale Inputs
% - kontinuierliche Change Streams
% - Safety-critical production use
%
% Diese Punkte in Outlook aufnehmen, nicht als aktuellen Scope darstellen.


% =============================================================================
\subsection{Beantwortung der Forschungsfragen}
\label{sec:discussion_rq_answers}
% =============================================================================

% WICHTIG:
% Die Forschungsfragen exakt aus Kapitel 1 uebernehmen.
% Nicht umformulieren.
%
% Hier spaeter jeweils:
%
% 1. exakte Forschungsfrage
% 2. kurze direkte Antwort (ca. 1 Absatz)
% 3. 3-5 tragende Ergebnisse
% 4. Literaturabgleich
% 5. Claim Boundary
%
% Reihenfolge:
% TF1 -> TF2 -> TF3 -> HF
%
% Dadurch baut die HF-Antwort auf den Teilfragen auf.


\subsubsection{Teilforschungsfrage 1: Input Assessment und Constraints}
\label{sec:answer_tf1}

% EXAKTE TF1:
%
% "Welche Anforderungen (Constraints) müssen an die Eingangsdaten gestellt werden,
% damit ein KI-Agent zuverlässig arbeiten kann?"
%
% MOEGLICHE ANTWORTSTRUKTUR, NOCH NICHT FINALE FORMULIERUNG:
%
% Constraints liegen auf mehreren Ebenen:
%
% A. Source Accessibility / Parsability
% B. klare Source Boundary
% C. ausreichende Evidence / Informationsgehalt
% D. Provenance / Identity
% E. Project Relevance / Project Fit
% F. Unsicherheit darf explizit bleiben
% G. keine ungepruefte Vermischung von Sources / Contexts
%
% EVIDENZ:
% SEM-001
% BLK-002
% BLK-003
% R4c
% Multi-Source Project 308131
%
% WICHTIG:
% "zuverlaessig" nicht als absolute fehlerfreie LLM-Leistung interpretieren.
% Definieren im Kontext der Architektur:
% reviewable, provenance-preserving, fail-closed bei Integritaetsproblemen.


\subsubsection{Teilforschungsfrage 2: Assurance und Effective Intervention}
\label{sec:answer_tf2}

% TF2 EXAKTEN WORTLAUT AUS KAPITEL 1 EINFUEGEN.
% Nicht aus Erinnerung paraphrasieren, sobald Kapitel 1 finalisiert ist.
%
% Inhaltliche Antwortstruktur:
%
% Human muss nicht jede automatisierte Operation kontrollieren.
%
% Human Authority ist erforderlich bei:
% - Engineering Meaning / Review
% - materieller Unsicherheit / Varianz
% - Model Placement
% - Formulation / Quality, wo semantische Entscheidung verbleibt
% - unresolved Relationships
% - Final Model Review
% - Publication
%
% Effective Intervention:
% - klare Evidence
% - sichtbare Alternatives / Variance
% - explicit Decision Semantics
% - Reopen / revise
% - keine substantive defaults
% - reduzierte Interaktion bei echtem Consensus
%
% EVIDENZ:
% SEM-005
% SEM-008
% SEM-013
% SEM-014
% OBS-020/021/025/026/031
% Human gates
%
% Safe-Demo-Feedback spaeter integrieren.


\subsubsection{Teilforschungsfrage 3: Leitplanken-Design}
\label{sec:answer_tf3}

% TF3 EXAKTEN WORTLAUT AUS KAPITEL 1 EINFUEGEN.
%
% Inhaltliche Antwortstruktur:
%
% Ontologische Regeln / Profiles sind notwendige,
% aber im PoC nicht hinreichende Leitplanken.
%
% Erforderliches Gesamtset:
%
% - controlled semantics / ontology
% - source grounding
% - provenance / identity
% - Project Fit
% - Human Authority
% - separation Meaning / Representation
% - explicit Placement
% - Target-Model Formulation
% - deterministic Assembly / Serialization
% - independent validation
%
% Plug & Play / Modularity:
% nicht durch Ontologie allein.
%
% Stattdessen:
% explizite Contracts und profile-controlled transformations.
%
% EVIDENZ:
% SEM-011
% SEM-012
% SEM-014
% SEM-015
% BLK-006
% BLK-007
% SYSIDE Validation
%
% Literatur:
% Ontology + transformation + SysML + provenance.


\subsubsection{Hauptforschungsfrage}
\label{sec:answer_hf}

% EXAKTE HF:
%
% "Wie muss eine Informationsverarbeitungsarchitektur gestaltet sein, um
% heterogene Bestandsdaten durch ein systematisches Bewertungs- und
% Synthese-Framework sicher in valide SysML v2 Strukturen zu überführen?"
%
% MOEGLICHE SYNTHETISCHE ANTWORTSTRUKTUR:
%
% Eine tragfaehige Architektur benoetigt eine kontrollierte Kette:
%
% Project / Source Context
% -> source-local Processing
% -> source-grounded Evidence
% -> shared Subject Identity
% -> Persona Interpretation
% -> Human Engineering Review
% -> Approved Engineering Information
% -> Project Fit / Multi-Source Reconciliation
% -> Human Project-level Authority
% -> Model Derivation / Placement
% -> Target-Model Formulation
% -> deterministic IEM Assembly
% -> deterministic SysML-v2 Generation
% -> independent Validation
% -> Final Human Model Review
% -> Publication Gate
%
% "sicher" in der Thesis nicht im funktional-safety-Sinn,
% sondern als governed / traceable / fail-closed / Human-authorized erklaeren,
% falls dieser Begriff in Kapitel 1 so gemeint ist.
%
% "valide SysML v2 Strukturen":
% formale / toolseitige Validitaet
% + separate Human Engineering Acceptance.
%
% WICHTIG:
% HF-Antwort muss kompakt sein.
% Keine erneute komplette Architekturkapitel-Zusammenfassung.


% =============================================================================
\subsection{Synthese der Diskussion}
\label{sec:discussion_synthesis}
% =============================================================================

% ZWECK:
% 1 Seite oder weniger.
%
% Zusammenfassen:
%
% - Welche Architekturannahmen wurden bestaetigt?
% - Welche mussten korrigiert / differenziert werden?
% - Was ist der eigentliche wissenschaftliche / konzeptionelle Beitrag?
% - Welche Claim Boundaries bleiben?
%
% MOEGLICHE VIER-PUNKT-SYNTHESE:
%
% 1. Grounding / Provenance vor freier KI-Interpretation.
% 2. Human Authority an materiellen Engineering-Entscheidungen.
% 3. Meaning / Representation / Formulation / Serialization trennen.
% 4. Multi-Source braucht project-wide identity + explicit authority.
%
% Plus:
% formale Validation != Engineering Correctness.
%
% UEBERGANG ZU KAPITEL 9:
% Fazit fasst Beitrag und Ausblick komprimiert zusammen.


% =============================================================================
% EMPFOHLENE ABBILDUNGEN / TABELLEN
% =============================================================================
%
% PRIORITAET 1
% Tab. 8.1 Ergebnis -> Interpretation -> Literatur -> RQ
%
% Ergebniscluster
% | Interpretation
% | Literaturbezug
% | RQ
% | Claim Boundary
%
% PRIORITAET 1
% Abb. 8.1 Synthese der Architekturprinzipien
%
% Grounding
% -> Human Authority
% -> Model Responsibility Separation
% -> Deterministic Generation / Validation
%
% Diese Abbildung darf interpretativ sein.
%
% PRIORITAET 1
% Tab. 8.2 Beantwortung der Forschungsfragen
%
% RQ
% | Kurzantwort
% | wichtigste Evidenz
% | Grenze
%
% PRIORITAET 2
% Tab. 8.3 Threats to Validity
%
% Threat
% | Auswirkung
% | Mitigation
% | Rest-Risiko
%
% PRIORITAET 2
% Tab. 8.4 AI im System vs AI im Entwicklungsprozess
%
% Rolle
% | Aufgabe
% | Human Authority
% | Risiko
% | Kontrolle
%
% OPTIONAL
% Effizienz-/Token-Tabelle, NUR wenn belastbare Daten vorliegen.


% =============================================================================
% LITERATURBEDARF KAPITEL 8
% =============================================================================
%
% HOCH:
%
% [ADD] Wieringa 2014
% [ADD] HITL / Human oversight:
%       Lazaros 2026
%       Mosqueira-Rey 2023
%       Zhu 2026
%
% [ADD] Multi-agent variance / disagreement:
%       Wang 2023
%       Liang 2024
%       Du 2024
%       Cui 2026
%       Kostka & Chudziak 2026
%
% [ADD] W3C PROV-DM / PROV Primer
%
% [ADD] Architecture as Code:
%       Bucaioni ICSA 2025
%       Stirbu IEEE Software 2022
%       optional Pontillo ICSA 2026
%
% [ADD] Topcu & Husain 2025 Trust at Your Own Peril
%
% [ADD] Prompt-Strategy-Driven SysML-v2 Artefact Generation 2026
%
% MITTEL:
%
% [ADD] AI coding assistant empirical literature
% fuer die Development Reflection.
%
% [REPO-BIB]
% SysMLv2 official OMG
% Cibrian 2025 / 2025b
% Pan 2025
% Chis 2025
% Tian 2025
% Meyers 2013
% Kapos 2014
% Li 2022
% Zhou 2024
% Zhao 2023
% ontology papers
% ISO42030, falls passender Claim
%
% WICHTIG:
% Jede allgemeine Aussage vor Schreiben gegen Volltext pruefen.


% =============================================================================
% PROJEKTEVIDENZ FUER KAPITEL 8
% =============================================================================
%
% Kapitel 8 sollte primaer auf Kapitel 7 referenzieren.
% Fuer Detailbelege:
%
% [PROJECT]
% wp12_findings.md
% wp12_formative_self_test_report.md
% BLK-002 closeout
% Gate-3 checkpoints
% final CATIA / thesis transition SSOT
% selected ADRs as Design Rationale
%
% ADRs insbesondere:
% ADR-016 Human authority
% ADR-019 IEM
% ADR-021 deterministic generation
% ADR-022 validation
% ADR-023 final review/publication
% ADR-027 Evidence before Persona
% ADR-029 placement / assembly separation
% ADR-032/033/034 Multi-Source
%
% ADRs NICHT als wissenschaftliche Stuetzung verwenden.


% =============================================================================
% OFFENE ARBEIT VOR AUSFORMULIERUNG
% =============================================================================
%
% PRIORITAET A
%
% [ ] Exakte TF2- und TF3-Wortlaute aus finalem Kapitel 1 einfuegen.
%
% [ ] Literatur-Jagdliste nach Abschluss aller Kapitelgerueste erzeugen.
%
% [ ] P1 Findings -> Interpretationscluster final mappen.
%
% [ ] Safe Demo durchfuehren und Kapitel 7 zuerst aktualisieren.
%
% [ ] Danach erst Safe-Demo-Interpretation in Kapitel 8.
%
% [ ] Token / Laufzeit / Kosten Daten pruefen:
%     Was ist tatsaechlich geloggt und thesis-verwendbar?
%
% PRIORITAET B
%
% [ ] Finding -> ADR Mapping fuer Architekturkorrekturen verifizieren.
%
% [ ] Ontology-Claims aus altem LSR reduzieren / korrigieren.
%
% [ ] Literature Gap gegen finalen LSR / PRISMA Corpus pruefen.
%
% [ ] RQ Answers erst nach Literaturintegration finalisieren.
%
% PRIORITAET C
%
% [ ] AI-assisted development reflection mit geeigneter Literatur absichern.
%
% [ ] Transferability Claims bewusst konservativ formulieren.


% =============================================================================
% CLAIM-CHECK VOR DEM SCHREIBEN
% =============================================================================
%
% [ ] Jede Interpretation hat ein Ergebnis aus Kapitel 7?
% [ ] General Claims wissenschaftlich belegt?
% [ ] Keine neue empirische Evidenz nur in Kapitel 8?
% [ ] Keine Finding-ID als Ersatz fuer Argumentation?
% [ ] Ontologie nicht als alleinige Hallucination Guardrail?
% [ ] Human Authority nicht als pauschal "mehr Human = sicherer" dargestellt?
% [ ] Effective Intervention beruecksichtigt Interaktionsaufwand?
% [ ] Meaning / Representation / Formulation / Serialization getrennt?
% [ ] SYSIDE-Validitaet klar von Engineering Correctness getrennt?
% [ ] Testanzahl nicht als wissenschaftliche Validitaet?
% [ ] Multi-Source nicht uebergeneralisiert?
% [ ] Target Coverage Partial klar genannt?
% [ ] Safe Demo nicht als representative Expert Validation?
% [ ] AI-as-system und AI-as-development-tool getrennt?
% [ ] AI coding productivity nicht ohne Daten generalisiert?
% [ ] TF1 / TF2 / TF3 exakt aus Kapitel 1 zitiert?
% [ ] HF exakt aus Kapitel 1 zitiert?
% [ ] RQ-Beantwortung direkt, knapp und evidenzbasiert?
% [ ] Keine Forschungsfrage nachtraeglich umgeschrieben?
% [ ] Limitationen nicht erst im Fazit versteckt?
% [ ] Keine Gedankenstriche als Satzzeichen im finalen Fliesstext?
% [ ] Kein unpersoenliches "man"?
% [ ] Tabellenueberschrift oben, Abbildungsunterschrift unten?
%
% =============================================================================
